\chapter{Resultados e Discussões}

Este capítulo apresenta a análise crítica dos resultados obtidos com a implantação do \textit{pipeline} em ambiente de produção. O objetivo é avaliar a eficácia da solução proposta frente aos problemas de despadronização e morosidade detalhados no Capítulo 2, dimensionando os impactos quantitativos de desempenho e os ganhos qualitativos de governança de dados para a Veeries.

\section{Impacto Operacional e Dimensionamento do Processamento}

O benefício imediato e mais expressivo da automatização foi o ganho de eficiência temporal em escala. No paradigma anterior, a equipe dedicava cerca de 3 horas de trabalho humano aplicadas à leitura, transcrição e validação não estruturada em cada lote mensal de relatórios. Com a nova arquitetura, o tempo de extração, validação, cruzamento e consolidação dos dados despencou para um intervalo de execução entre 2 a 5 segundos.

Para dimensionar o impacto quantitativo da solução, a consolidação da base histórica processou de forma automatizada documentos abrangendo o período de janeiro de 2023 a fevereiro de 2026. Devido ao padrão de republicação retroativa da Fonte X, o \textit{pipeline} ingere um volume massivo de documentos: são 144 arquivos anuais apenas para esta origem (divididos entre publicações mensais e acumuladas), o que resulta em picos de até 23 relatórios processados simultaneamente em um único mês. A Fonte Y, de publicação linear, agregou 29 documentos históricos ao escopo.

Essa extração em larga escala gerou uma base rastreável de 24.588 registros consolidados na Camada \textit{Silver} (L1), sendo 23.270 provenientes da Fonte X e 1.318 da Fonte Y. Na Camada \textit{Gold} (L2), após a operação estrutural de despivoteamento (\textit{Melt}) das variáveis originais para 20 colunas de análise contínuas, o sistema computou um crescimento exponencial, validando 492.540 registros históricos no formato longo (sendo 238.260 registros mensais e 254.280 acumulados). 

Por fim, o filtro lógico de deduplicação da Camada \textit{View} cumpriu o seu papel de abstrair a redundância gerada pela sobreposição de relatórios da fonte, entregando 29.322 instâncias definitivas (28.584 da Fonte X e 738 da Fonte Y). O esforço da equipe, portanto, deixou de ser a transcrição braçal de arquivos e passou a ser o consumo imediato de uma base cronologicamente atualizada de quase 30 mil linhas fidedignas prontas para as ferramentas de \textit{Business Intelligence}.

\section{Governança, Resiliência e Tratamento de Exceções}

Sob a perspectiva qualitativa, a adoção da Arquitetura \textit{Medallion} e da estratégia \textit{fail-fast} (falhar rapidamente) provou-se altamente eficaz na proteção do ambiente analítico da empresa. Durante os ciclos de operação real, o sistema de monitoramento conseguiu identificar e barrar a propagação de anomalias que, no método manual, passariam despercebidas.

Com o novo fluxo, quando o \textit{pipeline} reporta uma interrupção de sucesso, o trabalho do desenvolvedor responsável restringe-se a investigar a causa apontada pelos \textit{logs} e atuar em três frentes principais de exceção:

\begin{itemize}
    \item \textbf{Produtos não mapeados:} Como os relatórios externos frequentemente introduzem categorias de produtos voláteis ou nomenclaturas inéditas, a validação de dicionário da camada \textit{Silver} (L1) barra o processamento instantaneamente. O desenvolvedor então apenas adiciona a nova instância ao dicionário central de mapeamento e reexecuta o código.
    \item \textbf{Inconsistências de estrutura nativa:} Quando a fonte altera ligeiramente a forma que publica os dados, o desenvolvedor realiza ajustes pontuais na camada L0 Revisada. Isso valida a premissa arquitetural discutida no Capítulo 3: é possível corrigir o formato de leitura sem perder a rastreabilidade original do documento preservado na L0 Original.
    \item \textbf{Auditoria de séries históricas e erros da fonte:} O impacto mais notável de qualidade de dados ocorreu durante a consolidação histórica das bases. As validações matemáticas identificaram o aparecimento de uma categoria de produto que não possuía sentido lógico na série temporal. A restrição estrutural revelou que a falha não era do algoritmo de extração, mas sim um erro de publicação da própria fonte externa emissora do relatório. 
\end{itemize}

O bloqueio automático desse erro de publicação externa comprova o sucesso primário deste projeto. A implementação não apenas resolveu o problema de lentidão, mas instituiu uma malha fina de confiabilidade sistêmica. A Veeries agora possui a garantia de que os relatórios que abastecem suas ferramentas de tomada de decisão (BI) são estritamente fidedignos e auditáveis, eliminando o risco do consumo de dados corrompidos.